\SetPageTheme{story}
\PageHeader
  {Structura cat timp}
  {Repetitia apare atunci cand un proces trebuie sa continue pana cand starea jocului se schimba suficient.}
  {Cat timp 1 din 6}

\begin{missionbox}[Unde o intalnim]
Intr-un joc, structura \texttt{cat timp} apare natural in multe locuri:
\begin{itemize}
\item un cronometru merge pana la zero;
\item un efect de protectie ramane activ cat timp mai are durata;
\item valurile de inamici continua cat timp nu s-a atins limita;
\item o animatie se deruleaza cat timp nu s-a terminat traseul.
\end{itemize}

Ca sa o programezi bine, ai nevoie de trei piese: start, conditie de continuare si pas de actualizare.
\end{missionbox}

\BlocklyExperiment{cq-i1-while-counter}{Numaratoare controlata}{Leaga valoarea initiala, conditia si actualizarea, apoi observa de ce o bucla buna trebuie sa se apropie de oprire.}

\clearpage

\SetPageTheme{theory}
\PageHeader
  {Anatomia unei bucle}
  {Verific, execut, actualizez, ma intorc la verificare.}
  {Cat timp 2 din 6}

\begin{examplebox}[Modelul minim]
\begin{lstlisting}[style=codequest]
i <- 1
cat timp (i <= 5) executa
    afiseaza i
    i <- i + 1
sfarsit
\end{lstlisting}
\end{examplebox}

\begin{missionbox}[Cum o urmarim corect]
Nu sarim direct la raspunsul final. Urmarim bucla iteratie cu iteratie:
\begin{enumerate}
\item verificam daca regula este adevarata;
\item executam corpul;
\item modificam valoarea care controleaza bucla;
\item revenim la conditie.
\end{enumerate}
\end{missionbox}

\begin{guidebox}[Trace complet]
Completeaza tabelul:
\begin{center}
\begin{tabular}{llll}
\toprule
Iteratia & i la intrare & Ce se afiseaza & i dupa actualizare \\
\midrule
1 & & & \\
2 & & & \\
3 & & & \\
4 & & & \\
5 & & & \\
\bottomrule
\end{tabular}
\end{center}
\AnswerSpace{5}
\end{guidebox}

\clearpage

\SetPageTheme{guided}
\PageHeader
  {Cum se opreste}
  {O bucla sanatoasa nu doar porneste corect. Ea se si termina corect.}
  {Cat timp 3 din 6}

\begin{conceptbox}[Riscul principal]
Daca pasul nu duce variabila mai aproape de oprire, bucla poate continua la nesfarsit. Asta nu este doar o mica greseala de calcul. Intr-un joc real poate bloca logica, poate consuma resurse si poate face experienta inutilizabila.
\end{conceptbox}

\begin{semibox}[Bucla se opreste sau nu?]
Analizeaza fiecare situatie:
\begin{enumerate}
\item \texttt{i <- 1}, conditie \texttt{i <= 5}, pas \texttt{i <- i + 1}
\item \texttt{i <- 1}, conditie \texttt{i <= 5}, pas \texttt{i <- i - 1}
\item \texttt{timp <- 10}, conditie \texttt{timp > 0}, pas \texttt{timp <- timp - 2}
\item \texttt{energie <- 8}, conditie \texttt{energie >= 0}, pas \texttt{energie <- energie - 3}
\end{enumerate}
Pentru fiecare, spune daca bucla se opreste si de ce.
\AnswerSpace{8}
\end{semibox}

\clearpage

\SetPageTheme{guided}
\PageHeader
  {Cronometre si cooldown-uri}
  {Aici structura cat timp incepe sa semene cu adevarata logica de gameplay.}
  {Cat timp 4 din 6}

\begin{guidebox}[Trei probleme de joc]
Pentru fiecare situatie, scrie:
\begin{itemize}
\item valoarea initiala;
\item conditia de continuare;
\item actualizarea de la fiecare pas.
\end{itemize}

\begin{enumerate}
\item Un cronometru care porneste la 30 si ajunge la 0.
\item Un cooldown care porneste la 5 si scade pana cand abilitatea devine disponibila.
\item Un numarator de valuri care merge de la 1 la 8.
\end{enumerate}
\AnswerSpace{10}
\end{guidebox}

\begin{semibox}[Reflectie scurta]
Care dintre cele trei exemple ti se pare cel mai natural pentru \texttt{cat timp} si de ce?
\AnswerSpace{4}
\end{semibox}

\clearpage

\SetPageTheme{simulation}
\PageHeader
  {Simulam ritmul buclei}
  {Cand schimbi pasul sau conditia, vezi imediat diferenta dintre o repetitie eficienta si una stangace.}
  {Cat timp 5 din 6}

\SimulatorCard{cq-i1-while-counter}{Contor repetitiv}{Relatia dintre conditie, actualizare si numarul total de iteratii.}{Ruleaza aceeasi bucla cu pasi diferiti si compara cate iteratii apar, ce valori sunt afisate si unde se opreste procesul.}

\begin{semibox}[Dupa scanare]
\begin{enumerate}
\item Ce pas produce o parcurgere lenta, dar foarte precisa?
\item Ce pas produce o parcurgere rapida, dar poate sari valori importante?
\item Cum ai alege pasul pentru un cronometru si cum l-ai alege pentru deplasarea unui obiect pe ecran?
\end{enumerate}
\AnswerSpace{6}
\end{semibox}

\clearpage

\SetPageTheme{challenge}
\PageHeader
  {Buclele tale}
  {Acum scrii singur bucle complete si justifici de ce sunt sigure.}
  {Cat timp 6 din 6}

\SyntaxCheck{cat-timp}

\begin{solobox}[Set complet]
\begin{enumerate}
\item Scrie o bucla care afiseaza: 5, 10, 15, 20, 25.
\item Scrie o bucla descrescatoare pentru numaratoarea: 12, 9, 6, 3, 0.
\item Scrie o bucla pentru un cooldown care se opreste exact la zero, fara sa coboare sub zero.
\item Explica in 4--5 randuri ce face diferenta dintre o bucla care functioneaza corect si una care se poate bloca.
\end{enumerate}
\AnswerSpace{12}
\end{solobox}

\begin{checkpointbox}[Ideea care ramane]
Structura \texttt{cat timp} este forma cea mai clara pentru procese dependente de stare: timp ramas, energie ramasa, progres atins sau resurse inca disponibile.
\end{checkpointbox}

\clearpage
